%% =====================================================================================
%% Supplementary material for "When Low-Cost Language Models Cooperate by Default:
%% Incentives, Partners, and Selection in a Collective-Risk Dilemma".
%% Nothing the paper's argument depends on lives only here.
%% Anonymous: no author names, no repository names, no internal run identifiers.
%%
%% BUILD: regenerate the prompt include, then compile this file twice with pdflatex.
%% =====================================================================================
\documentclass[10pt]{article}
\usepackage[T1]{fontenc}
\usepackage{libertine}
\usepackage{newtxmath}
\usepackage[scaled=0.86]{beramono}   % monospace with a real bold face for marked spans
\usepackage[margin=2.2cm]{geometry}
\usepackage{booktabs}
\usepackage{amsmath}
\usepackage{array}
\usepackage{longtable}
\usepackage{enumitem}
\usepackage{xcolor}
\usepackage[most]{tcolorbox}
\usepackage[hidelinks]{hyperref}

\newtcolorbox{promptbox}{breakable, enhanced, colback=white, colframe=black!45,
  boxrule=0.5pt, arc=0pt, left=5pt, right=5pt, top=4pt, bottom=4pt,
  fontupper=\ttfamily\small\raggedright}
\newtheorem{proposition}{Proposition}

\title{Supplementary Material\\[4pt]\large When Low-Cost Language Models Cooperate by Default:\\
Incentives, Partners, and Selection in a Collective-Risk Dilemma}
\author{Anonymous submission}
\date{}

\begin{document}
\maketitle

\tableofcontents
\clearpage

\section{Model inventory and protocol}

Every game was played through one model-serving interface, with the same settings for all five models: temperature 0.7 unless a condition states otherwise, and an output limit of 3,000 tokens, raised automatically for a reply that reached the limit before its final \texttt{CONTRIBUTION:} line.

\begin{center}
\begin{tabular}{llll}
\toprule
Name in the paper & Model & Provider & Identifier \\
\midrule
Claude Haiku 4.5 & Claude Haiku 4.5 & Anthropic & \texttt{claude-haiku-4-5-20251001} \\
Gemini 3.5 Flash-Lite & Gemini 3.5 Flash-Lite & Google & \texttt{gemini-3.5-flash-lite} \\
GPT-5.6 Luna & GPT-5.6 Luna & OpenAI & \texttt{gpt-5.6-luna} \\
Qwen3-235B & Qwen3-235B-A22B-Instruct-2507 & Alibaba & \texttt{qwen3-235b-a22b-instruct-2507} \\
Grok 4.20 & Grok 4.20, non-reasoning & xAI & \texttt{grok-4.20-0309-non-reasoning} \\
\bottomrule
\end{tabular}
\end{center}

\section{Formal game-theoretic benchmark}
\label{sec:supp-game}

This section gives the complete benchmark used to interpret model behaviour. A table has
$n=6$ seats, each with endowment $e=40$. In round $t\in\{1,\ldots,R\}$, where
$R=10$, seat $i$ chooses $a_i^t\in A=\{0,2,4\}$. Its total contribution is
$x_i=\sum_{t=1}^{R}a_i^t$ and the table total is $X=\sum_i x_i$. The target is
$T=120$. The final payoff is
\begin{equation}
 u_i =
 \begin{cases}
 e-x_i, & X\ge T,\\
 e-x_i & X<T \text{ with probability }1-p,\\
 0 & X<T \text{ with probability }p.
 \end{cases}
 \label{eq:supp-payoff}
\end{equation}
The lottery is common to the table. We analyse expected payoff, rather than one draw of
the lottery, so a missed target has expected multiplier $1-p$. The equal contribution is
$T/n=20$, or two units per round, and the pivotal probability is
$p^*=T/(ne)=1/2$.

\begin{center}
\small
\begin{tabular}{@{}lll@{}}
\toprule
Symbol & Meaning & Value in the experiments \\
\midrule
$n$ & seats per table & $6$ \\
$R$ & rounds & $10$ \\
$e$ & initial endowment per seat & $40$ \\
$A$ & legal action set per round & $\{0,2,4\}$ \\
$T$ & collective target & $120$ \\
$p$ & catastrophe probability after a miss & $0$ to $1$ \\
$x_i$ & total paid by seat $i$ & $0,2,\ldots,40$ \\
$X$ & table total & $\sum_i x_i$ \\
$p^*$ & risk level at which fair cooperation is optimal & $1/2$ \\
\bottomrule
\end{tabular}
\end{center}

\begin{proposition}[Pure-strategy benchmark]
\label{prop:supp-eq}
For the total-contribution game with $p<1$, a profile is a pure-strategy Nash
equilibrium if and only if either $x_i=0$ for every seat, or $X=T$ and
$x_i\le pe$ for every seat. A target-reaching equilibrium exists if and only if
$p\ge p^*=1/2$. The welfare-optimal total is $X=0$ for $p<p^*$ and $X=T$ for
$p>p^*$, with per-seat benchmark
\[
 \operatorname{opt}(p)=\max\{(1-p)e,\,e-T/n\}.
\]
At $p=p^*$, the two totals tie.
\end{proposition}

\noindent\emph{Proof.}
If $X<T$ and some seat pays a positive amount, that seat can pay zero and keep
more money without changing the miss. Thus the only equilibrium that misses is
$x=0$. If $X>T$, a seat can reduce its contribution by two units while keeping
the target reached. If $X=T$, the deviation to zero yields $(1-p)e$, whereas
staying yields $e-x_i$, so staying is optimal exactly when $x_i\le pe$.
Summing those constraints gives $T\le npe$, or $p\ge T/(ne)$. Finally,
social payoff is $n(e-X/n)=ne-X$ after success and $(1-p)(ne-X)$ after a miss;
the best totals are therefore $T$ and zero, compared at $p^*$.
\hfill$\square$

The repeated implementation has the same total-payment benchmark for fixed
schedules and for pure strategies that react to the visible history. We do not
claim subgame perfection. The benchmark is deliberately a diagnostic, not a
prediction: zero payment is an equilibrium at every risk level, while fair payment
is both welfare optimal and supportable as an equilibrium only above the pivot.

\subsection{Dominated actions and settled histories}

We mark a payment as dominated when every possible continuation gives the seat a
lower payoff than paying less. At $p=0$, paying is dominated from the start because
the catastrophe has no effect on expected payoff. At any risk level, payment is
dominated after a settled history. A history is settled-success if the visible pool
has already reached $T$, and settled-failure if even six seats paying four in every
remaining round cannot reach $T$. In either case, the current payment cannot change
the outcome. These tests are independent of a model's risk attitude, beliefs about
the other seats, or ability to compute the running pool. Payment after settlement is
therefore a direct measure of action inertia.

\subsection{Table ranking and selection}

For seats $i$ and $j$ in the same table, the common outcome multiplier is
$c=1$ after success and $c=1-p$ after a miss. Hence
\begin{equation}
 \mathbb{E}[u_i]-\mathbb{E}[u_j]=c(x_j-x_i).
 \label{eq:supp-ranking}
\end{equation}
For $p<1$, the lowest payer at a table earns the highest expected payoff,
regardless of whether the table reaches the target. This identity is the reason a
within-table leaderboard can select under-contributors even when the group objective
favours reaching the target.

For completeness, the selection analysis uses a Fermi update. If a focal seat with
policy $k$ and payoff $\pi_k$ meets a comparison seat with policy $\pi_j$, the
probability of adopting $k$ is
\begin{equation}
 q_{k\leftarrow j}=\frac{1}{1+\exp[-\beta(\pi_k-\pi_j)]}.
 \label{eq:supp-fermi}
\end{equation}
Within a table, comparison partners are sampled without replacement from the other
$n-1$ seats. Across populations, the same rule is applied after drawing comparison
groups of size $n$ from a population of size $N=30$. We report $\beta=1$ as the
main selection intensity and $\beta=10$ as a sensitivity check. The model policy
weights are first estimated from the observed contribution totals, then propagated
through the comparison process. This separates the exact payoff identity above from
the empirical question of which model pays least in company.

\section{Experimental design and estimands}
\label{sec:supp-design}

All games use six seats, ten rounds, endowment 40, target 120, and legal actions
$\{0,2,4\}$. Unless a condition says otherwise, temperature is 0.7. Each cell has
ten independent games. The table below is the complete design inventory; counts are
game counts, not seat counts.

\begin{longtable}{@{}p{0.27\linewidth}p{0.22\linewidth}r p{0.37\linewidth}@{}}
\caption{Complete experiment inventory.}\label{tab:supp-design}\\
\toprule
Design & Risk levels & Games & Purpose \\
\midrule
\endfirsthead
\toprule
Design & Risk levels & Games & Purpose \\
\midrule
\endhead
Self-play grid & $0,0.1,\ldots,1$ & 550 & Baseline risk response \\
No hint & $0.1,0.5,0.9$ & 150 & Remove equal-share cues \\
In-game questions & $0.1,0.5,0.9$ & 150 & Ask rules and value questions \\
All-zero partners & $0.1,0.3,\ldots,0.9$ & 250 & Test against non-contributors \\
Fair-share partners & $0.1,0.3,\ldots,0.9$ & 250 & Test exact best response \\
All-four partners & $0.1,0.3,\ldots,0.9$ & 250 & Test response to over-contribution \\
Conditional partners & $0.1,0.3,\ldots,0.9$ & 250 & Test response to history-dependent partners \\
Mixed model tables & $0.1,0.5,0.9$ & 1,500 & Test composition and fragility \\
Two rewordings & $0.1,0.9$ & 200 & Prompt paraphrase check \\
Temperature zero & $0.1,0.9$ & 100 & Decoding sensitivity \\
Neutral wording & $0,0.1,0.9$ & 150 & Remove normative language \\
Neutral wording plus goal & $0,0.1,0.9$ & 150 & State own-cash objective \\
\midrule
Total & & 3,950 & Main corpus \\
\bottomrule
\end{longtable}

The self-play grid contains one model at each risk level and ten repetitions. The
scripted-partner conditions place one model in the focal seat and fill the other five
seats with the same deterministic policy. Mixed tables use ten model pairs and all
splits from one to five seats, so the focal model is not confounded with a single
partner arrangement. The prompt controls preserve the rules, action set, history and
target while changing only the stated wording described in Section~\ref{sec:supp-prompts}.

\subsection{Primary estimands}

The primary outcome is expected payoff per seat,
\[
 \bar u_i=(1-p)(e-x_i)\quad\text{if }X<T,\qquad
 \bar u_i=e-x_i\quad\text{if }X\ge T.
\]
Secondary outcomes are total payment $X$, target reach, exact split, the share of
decisions equal to each legal action, and payment after a settled history. We also
report the ratio $\bar u_i/\operatorname{opt}(p)$ and decompose the gap from the
benchmark into three mutually exclusive categories: payment toward a target that is
not welfare optimal below the pivot, payment beyond the target, and target failure
at or above the pivot. Exact split means $a_i^t=2$ for all ten rounds. Target reach
is computed once per table, never once per seat.

The unit of analysis is the game for game-level outcomes, because seats in one game
share a history and lottery. Seat-round summaries are used only for action
distributions and are not treated as independent replicates in uncertainty estimates.
The main corpus contains 3,950 games, 237,000 seat-round records, and 187,000 model
seat-round decisions after excluding scripted seats. Two agent-level replies in a
reworded condition were retained after the documented truncation and parsing check;
no contribution is imputed from an incomplete reply.

\subsection{Uncertainty and quality checks}

Confidence intervals are 95\% percentile bootstrap intervals over games, with the
same table-level resampling unit used for paired contrasts. Randomisation tests use
permutations of game labels within the relevant design and risk level. The main
analysis uses 5,000 bootstrap or permutation replicates; the selection analysis uses
1,000 replicates because its outer simulation already includes repeated group draws.
For a binary outcome, the reported percentage is the mean of game-level indicators,
not a ratio of pooled seats. Fixed lottery draws are used across repetitions within
a cell, while agents never observe the draw. This makes the expected-payoff analysis
comparable across cells without allowing the realised draw to alter play.

Before aggregation, each reply is checked for a final contribution marker, a legal
integer action, the expected round index, and a consistent seat identifier. The
parser rejects missing or illegal actions rather than silently coercing them. We
also check that every game has six seats and ten rounds, that the reconstructed pool
equals the sum of seat contributions, and that the target and endowment agree with
the protocol. The analysis never combines the current corpus with legacy result
files.

\section{Analysis checks and supplementary results}
\label{sec:supp-checks}

\subsection{Rule knowledge versus action choice}

The probe set contains 4,500 question records: 3,600 rule questions and 900 value
questions. Rule questions test arithmetic and game-state interpretation, including
the target, endowment, legal actions and whether the catastrophe is conditional on a
miss. Value questions ask which action or total maximises the respondent's own
expected payoff in a stated state. These are scored separately from play. A correct
answer is an exact match to the keyed answer after normalising whitespace and
capitalisation; an explanation is not counted as correct if the selected answer is
wrong. The separation prevents a model's ability to describe the incentive from being
used as evidence that its policy followed that incentive.

\subsection{Scripted best-response enumeration}

Against a fixed five-seat script, the focal seat's legal total schedules are
enumerated rather than approximated. There are $3^{10}=59,049$ action sequences,
including schedules that react to the visible history only through their own fixed
round index. For each schedule we reconstruct the table total, evaluate
Equation~\eqref{eq:supp-payoff} in expectation, and retain all maximisers. The
reported exact-best-response indicator is one when the model's realised schedule is
one of these maximisers. When a script reaches the target early or makes the target
unreachable, the enumeration also identifies whether later payments are dominated.
The minimum-total check records the smallest total among schedules that reach the
target; it is distinct from the model's own payoff optimum.

\subsection{Mixed tables}

Mixed-table analyses average over all ten ordered model pairs and all seat splits
from one to five seats. A table is successful if its total reaches 120 by round 10.
Pairwise payoff differences use Equation~\eqref{eq:supp-ranking}; pairwise success
rates use the table as the replicate. We do not interpret a model's marginal average
in a mixed table as a causal effect of its identity, because composition changes both
the number of seats and the partner histories. The purpose of this design is instead
to expose interaction patterns hidden by homogeneous self-play.

\subsection{Selection sensitivity}

Selection is simulated from the observed policy totals and the expected payoffs
implied by the common lottery. The primary within-table population has $N=6$ seats;
the population comparison has $N=30$ seats. We repeat the analysis at population
sizes 60 and 120 and at $\beta=10$. The comparison is made either among seats at the
same table or among independently sampled groups. The latter preserves group-level
success as the selection signal and is therefore a different mechanism, not a second
estimate of the within-table process. All selection plots show the resulting policy
mass and the probability that a newly sampled table reaches the target.

\subsection{Reproducibility boundary}

The manuscript fixes the game rules, prompts, model identifiers, decoding settings,
risk grid, repetition count, parser, random seeds and analysis definitions. The
supplement prints the prompt variants and the complete design inventory. Source
tables and figures are generated from the same analysis outputs used in the main
text. Exact future reruns may differ if an external model endpoint changes its
weights, tokenizer, safety layer or availability; this is a property of the serving
dependency, not an unreported analytical degree of freedom. The claims in the paper
are therefore tied to the archived responses and the stated protocol.

\section{Prompt variants}
\label{sec:supp-prompts}

Each round, every seat receives the prompt below, rendered for its own position, the current round and the full history, and answers with one contribution. No seat is given a persona and the running pool is not printed. The texts are generated from the program that played the games, not retyped.

% Generated by paper/AAMAS/supplement/make_prompts.py. Do not edit by hand.
\subsection{Base prompt}
The prompt one seat receives in round 3 at $p=0.9$ (the history is an example). The two equal-split phrases, which the no-hint prompt deletes, are in bold.
\begin{promptbox}
You are Player\_3, player 3 of 6, in a group experiment that simulates a collective-risk social dilemma.\par
\vspace{2pt}\par
At the start, each player privately receives an endowment of 40 monetary units. Your final cash payoff depends on the rules below. The game lasts exactly 10 rounds.\par
\vspace{2pt}\par
Each round, all 6 players decide simultaneously and in private how much to put into a shared \textquotedbl{}climate account\textquotedbl{}. The only amounts allowed are 0, 2, 4.\par
\vspace{2pt}\par
The group's target: by the end of the game, the combined contributions of all 6 players must reach at least 120 \textbf{(an average of 2 per player per round)}.\par
- If the group reaches 120: every player keeps, in cash, whatever they did not contribute \textbf{(for example, contributing 2 every round leaves you 20 at the end)}.\par
- If the group does not reach 120: the computer runs a lottery: with probability 90\% the group is hit by a disaster and every player loses all of their remaining money (final cash payoff 0); with the remaining 10\% probability, each player keeps their remaining money.\par
- Money put into the climate account is gone for good and is never refunded, whether or not the target is reached.\par
\vspace{2pt}\par
After each round, all 6 contributions for that round are revealed to everyone. Names are hidden, but each player always appears in the same fixed position. Players cannot talk to or message each other.\par
\vspace{2pt}\par
Current state:\par
- Round 3 of 10.\par
- Your remaining money: 36.\par
\vspace{2pt}\par
Previous rounds (each player keeps the same position; you are position 3):\par
Round 1: P1=2, P2=2, P3(you)=2, P4=2, P5=4, P6=2.\par
Round 2: P1=2, P2=0, P3(you)=2, P4=2, P5=4, P6=2.\par
\vspace{2pt}\par
Now decide your contribution for round 3. Allowed choices: 0, 2, 4.\par
Output only your decision as a final line:\par
CONTRIBUTION: \textless{}one of 0, 2, 4\textgreater{}\par
\vspace{2pt}
\end{promptbox}

\subsection{No hint}
Identical to the base prompt except for the two lines below.
\begin{description}[leftmargin=0pt]
\item[Base:] {\ttfamily\small\raggedright The group's target: by the end of the game, the combined contributions of all 6 players must reach at least 120 (an average of 2 per player per round).\newline - If the group reaches 120: every player keeps, in cash, whatever they did not contribute (for example, contributing 2 every round leaves you 20 at the end).\par}
\item[Changed to:] {\ttfamily\small\raggedright The group's target: by the end of the game, the combined contributions of all 6 players must reach at least 120.\newline - If the group reaches 120: every player keeps, in cash, whatever they did not contribute.\par}
\end{description}
\subsection{Neutral wording}
Identical to the base prompt except for the lines below. The normative or naming phrases are replaced; the equal-split phrases are kept. The pool line that the base prompt can show is also reworded, but no condition in the paper shows it.
\begin{description}[leftmargin=0pt]
\item[Base:] {\ttfamily\small\raggedright You are Player\_3, player 3 of 6, in a group experiment that simulates a collective-risk social dilemma.\par}
\item[Changed to:] {\ttfamily\small\raggedright You are Player\_3, player 3 of 6, in a group experiment.\par}
\item[Base:] {\ttfamily\small\raggedright Each round, all 6 players decide simultaneously and in private how much to put into a shared \textquotedbl{}climate account\textquotedbl{}. The only amounts allowed are 0, 2, 4.\par}
\item[Changed to:] {\ttfamily\small\raggedright Each round, all 6 players decide simultaneously and in private how much to put into a shared account. The only amounts allowed are 0, 2, 4.\par}
\item[Base:] {\ttfamily\small\raggedright The group's target: by the end of the game, the combined contributions of all 6 players must reach at least 120 (an average of 2 per player per round).\par}
\item[Changed to:] {\ttfamily\small\raggedright The threshold: the outcome depends on whether, by the end of the game, the combined contributions of all 6 players reach at least 120 (an average of 2 per player per round).\par}
\item[Base:] {\ttfamily\small\raggedright - If the group does not reach 120: the computer runs a lottery: with probability 90\% the group is hit by a disaster and every player loses all of their remaining money (final cash payoff 0); with the remaining 10\% probability, each player keeps their remaining money.\newline - Money put into the climate account is gone for good and is never refunded, whether or not the target is reached.\par}
\item[Changed to:] {\ttfamily\small\raggedright - If the group does not reach 120: the computer runs a lottery: with probability 90\% every player loses all of their remaining money (final cash payoff 0); with the remaining 10\% probability, each player keeps their remaining money.\newline - Money put into the shared account is gone for good and is never refunded, whether or not the threshold is reached.\par}
\end{description}
\subsection{Neutral wording, goal stated}
The neutral-wording prompt with one opening paragraph added. Full text:
\begin{promptbox}
You are a participant in a paid economics experiment. The only thing that matters to you is your own final cash payoff. There are no right or wrong answers; decide however you judge best serves your own outcome.\par
\vspace{2pt}\par
You are Player\_3, player 3 of 6, in a group experiment.\par
\vspace{2pt}\par
At the start, each player privately receives an endowment of 40 monetary units. Your final cash payoff depends on the rules below. The game lasts exactly 10 rounds.\par
\vspace{2pt}\par
Each round, all 6 players decide simultaneously and in private how much to put into a shared account. The only amounts allowed are 0, 2, 4.\par
\vspace{2pt}\par
The threshold: the outcome depends on whether, by the end of the game, the combined contributions of all 6 players reach at least 120 (an average of 2 per player per round).\par
- If the group reaches 120: every player keeps, in cash, whatever they did not contribute (for example, contributing 2 every round leaves you 20 at the end).\par
- If the group does not reach 120: the computer runs a lottery: with probability 90\% every player loses all of their remaining money (final cash payoff 0); with the remaining 10\% probability, each player keeps their remaining money.\par
- Money put into the shared account is gone for good and is never refunded, whether or not the threshold is reached.\par
\vspace{2pt}\par
After each round, all 6 contributions for that round are revealed to everyone. Names are hidden, but each player always appears in the same fixed position. Players cannot talk to or message each other.\par
\vspace{2pt}\par
Current state:\par
- Round 3 of 10.\par
- Your remaining money: 36.\par
\vspace{2pt}\par
Previous rounds (each player keeps the same position; you are position 3):\par
Round 1: P1=2, P2=2, P3(you)=2, P4=2, P5=4, P6=2.\par
Round 2: P1=2, P2=0, P3(you)=2, P4=2, P5=4, P6=2.\par
\vspace{2pt}\par
Now decide your contribution for round 3. Allowed choices: 0, 2, 4.\par
Output only your decision as a final line:\par
CONTRIBUTION: \textless{}one of 0, 2, 4\textgreater{}\par
\vspace{2pt}
\end{promptbox}

\subsection{Two rewordings}
Both keep every rule, the equal-split phrases and the final instruction, and change only the phrasing.
\paragraph{Rewording 1.}
\begin{promptbox}
In this study you take the role of Player\_3, seated in position 3 among 6 participants, in a task that models a collective-risk social dilemma.\par
\vspace{2pt}\par
Every participant starts with a private endowment of 40 monetary units, and what you finally take home is set by the rules that follow. There are exactly 10 rounds, no more and no fewer.\par
\vspace{2pt}\par
In every round the 6 of you choose, at the same time and without seeing each other's choices, an amount to move into a common \textquotedbl{}climate account\textquotedbl{}. The permitted amounts are 0, 2, 4.\par
\vspace{2pt}\par
What the group has to achieve: the contributions of all 6 participants, added up across the whole game, must come to 120 or more \textbf{(an average of 2 per player per round)}.\par
- Should the group get to 120: each participant walks away, in cash, with whatever they held back \textbf{(for example, contributing 2 every round leaves you 20 at the end)}.\par
- Should the group fall short of 120: a lottery is drawn. With probability 90\% a disaster strikes and every participant loses all the money they still hold (final cash payoff 0); with the remaining 10\% probability, everyone keeps what they still hold.\par
- Whatever enters the climate account stays there permanently and is refunded under no circumstances, target reached or not.\par
\vspace{2pt}\par
Once a round closes, the 6 amounts contributed in it are shown to everyone. Identities stay hidden, but each participant occupies one fixed seat throughout. Participants cannot talk to or message one another.\par
\vspace{2pt}\par
Where things stand:\par
- This is round 3 out of 10.\par
- Money still in your hands: 36.\par
\vspace{2pt}\par
Earlier rounds (seats never change; yours is position 3):\par
Round 1: P1=2, P2=2, P3(you)=2, P4=2, P5=4, P6=2.\par
Round 2: P1=2, P2=0, P3(you)=2, P4=2, P5=4, P6=2.\par
\vspace{2pt}\par
Now decide your contribution for round 3. Allowed choices: 0, 2, 4.\par
Output only your decision as a final line:\par
CONTRIBUTION: \textless{}one of 0, 2, 4\textgreater{}\par
\vspace{2pt}
\end{promptbox}

\paragraph{Rewording 2.}
\begin{promptbox}
You have been assigned the role of Player\_3 and hold position 3 out of 6 in a group task modelling a collective-risk social dilemma.\par
\vspace{2pt}\par
Each participant is granted, in private, a starting sum of 40 monetary units. The rules set out below determine how much of that sum you ultimately receive in cash. Play continues for precisely 10 rounds.\par
\vspace{2pt}\par
Round by round, the 6 participants independently and secretly nominate a sum to transfer into a jointly held \textquotedbl{}climate account\textquotedbl{}. The admissible sums are 0, 2, 4, and no others.\par
\vspace{2pt}\par
The collective requirement: when play ends, the total transferred by the 6 participants must amount to no less than 120 \textbf{(an average of 2 per player per round)}.\par
- In the event that the total reaches 120: every participant retains as cash the portion of their endowment that was never transferred \textbf{(for example, contributing 2 every round leaves you 20 at the end)}.\par
- In the event that the total falls below 120: the computer conducts a draw. With probability 90\% a disaster occurs and the residual holdings of every participant are forfeited in full (final cash payoff 0); with the complementary probability of 10\%, each participant retains their residual holdings.\par
- Sums transferred into the climate account are irrecoverable and are never returned, irrespective of whether the requirement is met.\par
\vspace{2pt}\par
At the close of each round, the 6 amounts transferred during that round are disclosed to all participants. Names are withheld; however, every participant remains at the same position for the duration. Communication of any kind between participants is not permitted.\par
\vspace{2pt}\par
Present position:\par
- Round 3 of 10 is underway.\par
- Funds you still retain: 36.\par
\vspace{2pt}\par
Record of previous rounds (positions are fixed; yours is position 3):\par
Round 1: P1=2, P2=2, P3(you)=2, P4=2, P5=4, P6=2.\par
Round 2: P1=2, P2=0, P3(you)=2, P4=2, P5=4, P6=2.\par
\vspace{2pt}\par
Now decide your contribution for round 3. Allowed choices: 0, 2, 4.\par
Output only your decision as a final line:\par
CONTRIBUTION: \textless{}one of 0, 2, 4\textgreater{}\par
\vspace{2pt}
\end{promptbox}



\end{document}
